% New author proof. Main inputs: original invariant collision system and native330.
% No exponential return tail or strong spectral theorem is assumed below.
\section{The actual return tower and positive finite-band likelihoods}
\label{sec:av-return-likelihood}

We retain the probability $\nu_R^*=\nu|_{Y_R^*}/c_*$ and the
unit-speed physical clock.  All maps in this section act on the
full-measure invariant set of trajectories whose forward and backward
collisions are regular.  This removes no positive collision mass and
assigns no value to a separate singular boundary trace.

\begin{lemma}[The conservative physical return tower]
\label{lem:av-tower}
Let $r_R(x)=\min\{j\ge1:T_R^jx\in Y_R^*\}$ be the true first
return in collisions and $F_R^*x=T_R^{r_R(x)}x$.  The return is finite
almost everywhere on the original section, $F_R^*$ preserves
$\nu_R^*$, and
\begin{equation}\label{eq:av-tower-moment}
 \int_{Y_R^*}r_R\,d\nu\le1,\qquad
 \mathbb E_{\nu_R^*}N_{n,R}\le\frac{n}{c_*},\qquad
 \nu_R^*\{N_{n,R}>L\}\le\frac{n}{c_*L}\quad(L>0),
\end{equation}
where
\[
 N_{n,R}=\sum_{a=0}^{n-1}r_R\circ(F_R^*)^a
\]
is the actual collision count of the $n$th return.  Its physical time
$S_{n,R}=L_{N_{n,R},R}$ satisfies
$\mathbb E_{\nu_R^*}S_{n,R}\le\tau_+n/c_*$.
\end{lemma}
\begin{proof}
Suppress $R$ and write $Y=Y_R^*$.  The actual invertible collision
map preserves $\nu$.  Let $E\subset Y$ consist of points with no
future visit to $Y$.  The sets $T^jE$, $j\ge0$, are disjoint: an
intersection at times $i<j$ would give a future visit to $Y$ at time
$j-i$ from a point of $E$.  Their equal measures have sum at most
one, so $\nu(E)=0$.  Points of $Y$ with only finitely many future
visits belong to $\bigcup_{j\ge0}T^{-j}E$, also null.  The same
argument for $T^{-1}$ gives infinitely many past visits almost
everywhere on $Y$.

Put $Y_a=\{x\in Y:r(x)=a\}$.  The sets $T^aY_a$, $a\ge1$,
partition $Y$ up to a null set: for almost every terminal point,
its nearest strictly earlier visit gives the unique starting point
and the unique positive return time.  Thus for every measurable
$B\subset Y$, invariance and disjointness give
\[
 \nu((F^*)^{-1}B)=\sum_{a\ge1}\nu(Y_a\cap T^{-a}B)
 =\sum_{a\ge1}\nu(T^aY_a\cap B)=\nu(B).
\]
Division by the original $c_*$ proves induced invariance.

The true tower levels $T^j\{x\in Y:r(x)>j\}$, $j\ge0$, are
disjoint.  Indeed, an intersection at levels $i<j$ implies that a
point with return time exceeding $j$ has a visit to $Y$ at time
$j-i\le j$, a contradiction.  Consequently
\[
 \int_Yr\,d\nu=\sum_{j\ge0}\nu\{x\in Y:r(x)>j\}
 =\nu\!\left(\bigsqcup_{j\ge0}T^j\{x\in Y:r(x)>j\}\right)
 \le1.
\]
Equality is not needed and is not asserted without saturation of
the collision space.  Induced invariance gives the expectation of
the sum defining $N_n$ and hence the second bound in
\eqref{eq:av-tower-moment}.  Integrating
$\one_{\{N_n>L\}}\le N_n/L$ proves its last bound.
Finally $S_n\le\tau_+N_n$ on each true trajectory.  No mixing,
independence, spectral estimate or ergodicity assertion enters this
argument.
\end{proof}

Let $K_{n,R}=K^{\rm c}_{N_{n,R},R}$ be the actual two-dimensional
lattice displacement at that return.  The full probability law is
\begin{equation}\label{eq:av-return-law}
 \mu_{n,R}=\operatorname{Law}_{\nu_R^*}
                   (K_{n,R},N_{n,R},S_{n,R})
       \quad\hbox{on }\mathbb Z^2\times\mathbb N\times\mathbb R.
\end{equation}
Fixing $n$ in this law retains the two displacement coordinates, true
collision count and physical time.  Fixing $m$ instead gives the
original coordinates $(k_1,k_2,t,n)$ in \eqref{eq:v76-target}.
These are two slices of the same joint records, not different
conditioning events or a replacement inverse problem.

\begin{lemma}[Exact components and a positive proof restriction]
\label{lem:av-return-components}
The $(k,m)$ component of \eqref{eq:av-return-law} is exactly
$p_{n,R}(k,m,t)\,dt$.  For an integer $L\ge1$ put
\[
 Q_{n,R}^{L}=\one_{\{m\le L\}}\mu_{n,R},\qquad
 H_{n,R}^{L}=\mu_{n,R}-Q_{n,R}^{L},\qquad
 \delta_{n,R}(L)=\|H_{n,R}^{L}\|_{\rm TV}.
\]
Both are positive measures, and
\begin{equation}\label{eq:av-return-bv}
 \delta_{n,R}(L)\le\frac{n}{c_*L},\qquad
 \sum_{m\le L}\sum_k\|D p_{n,R}(k,m,\cdot)\|_{\rm TV}
                  \le D_{\rm all}^{1}L.
\end{equation}
Neither measure is renormalized.  This restriction is a device of
the proof and is removed from the conclusions about $\mu_{n,R}$.
\end{lemma}
\begin{proof}
Starting in $Y_R^*$, the event $N_{n,R}=m$ says that the time-$m$
state is in that section and precisely $n$ visits, including the
initial one, occur at indices $0,\ldots,m-1$.  It is therefore
exactly the occupation and endpoint event in \eqref{eq:v76-event}.
Adding the true displacement label gives its stated density.  The
tail bound is \eqref{eq:av-tower-moment}.  At each $m$, fixing one
occupation label in Theorem~\ref{thm:av-all-bv} decreases its positive
variation sum.  Summing over at most $L$ actual collision counts
proves \eqref{eq:av-return-bv}.  Almost-sure successive recurrence
from Lemma~\ref{lem:av-tower} shows that all components together
have mass one.
\end{proof}

\subsection{A positive kernel with an exact compact frequency support}

\begin{lemma}[The kernel and its first absolute moment]
\label{lem:av-kernel}
With its continuous value at zero, put
\begin{equation}\label{eq:av-kernel}
 k(t)=\frac{3}{8\pi}\left(\frac{\sin(t/4)}{t/4}\right)^4,
 \qquad k_{\mathfrak b}(t)=\mathfrak b k(\mathfrak b t)
 \quad(\mathfrak b>0).
\end{equation}
Then $k\ge0$, $\int k=1$, $\int t^2k(t)\,dt=12$, and
$\kappa_{\rm ker}:=\int|t|k(t)\,dt\le\sqrt{12}$.
For the convention $\widehat f(s)=\int e^{ist}f(t)\,dt$,
\begin{equation}\label{eq:av-kernel-transform}
 \widehat k(s)=\eta(s)=
 \begin{cases}
 1-6s^2+6|s|^3,&|s|\le1/2,\\
 2(1-|s|)^3,&1/2\le|s|\le1,\\
 0,&|s|\ge1.
 \end{cases}
\end{equation}
The bandwidth $\mathfrak b$ is independent of the outer geometric
parameter $B$ in the original shortage limit.
\end{lemma}
\begin{proof}
The displayed $\eta$ is compactly supported and $C^2$.  Its third
derivative has jumps $12,-48,72,-48,12$ at respectively
$-1,-1/2,0,1/2,1$.  Four integrations by parts, or the same
distributional identity tested against $e^{-its}$, give for $t\ne0$
\[
 t^4\int\eta(s)e^{-its}\,ds
 =24\cos t-96\cos(t/2)+72=192\sin^4(t/4).
\]
Thus $(2\pi)^{-1}\int\eta(s)e^{-its}\,ds$ equals the $k$ in
\eqref{eq:av-kernel}; at zero both equal $3/(8\pi)$ since
$\int\eta=3/4$.  This also proves nonnegativity, and the displayed
formula gives an integrable $O(|t|^{-4})$ tail.  Fourier inversion
for this integrable transform pair gives $\int k=\eta(0)=1$.
Since $t^2k$ is integrable, differentiation of its transform is
justified by domination, and
$\int t^2k(t)\,dt=-\eta''(0)=12$.  Cauchy--Schwarz against the
probability density $k$ gives the asserted first absolute moment.
Scaling gives transform $\eta(s/\mathfrak b)$ and first absolute
moment $\kappa_{\rm ker}/\mathfrak b$ for $k_{\mathfrak b}$.
\end{proof}

Convolution below acts only on the real roof coordinate and preserves
every discrete label.  The norm $\|\cdot\|_{\rm TV}$ of a signed
measure is its full variation, without a factor one half.

\begin{theorem}[Positive likelihoods on the original trajectories]
\label{thm:av-likelihood}
Fix $n\ge1$, $L\ge1$, $\mathfrak b>0$, $R\in[R_-,R_+]$ and an
integer $J\ge1$.  For every label $\ell=(k,m)$ let $I_\ell$ be
empty or a union of at most $J$ bounded physical-roof intervals.
These intervals and labels specify the unchanged event
\[
 A=\{x:S_{n,R}(x)\in I_{(K_{n,R}(x),N_{n,R}(x))}\}.
\]
On the very same trajectories define the analytical likelihood
\begin{equation}\label{eq:av-likelihood-definition}
 \lambda_{A,\mathfrak b}(x)=
 (k_{\mathfrak b}*\one_{I_{(K_{n,R}(x),N_{n,R}(x))}})
                                            (S_{n,R}(x)).
\end{equation}
Then $0\le\lambda_{A,\mathfrak b}\le1$.  If
$\mathcal V_L=D_{\rm all}^{1}L$ and
$\delta=\delta_{n,R}(L)$, then
\begin{align}
 \|\mu_{n,R}-k_{\mathfrak b}*\mu_{n,R}\|_{\rm TV}
       &\le2\delta+\frac{\kappa_{\rm ker}\mathcal V_L}{\mathfrak b},
                                                       \label{eq:av-law-error}\\
 \mathbb E_{\nu_R^*}|\one_A-\lambda_{A,\mathfrak b}|
       &\le\delta+\frac{\kappa_{\rm ker}J\mathcal V_L}{\mathfrak b}
          =:\mathcal E_n(L,\mathfrak b,J),              \label{eq:av-event-error}\\
 \left|\mathbb E_{\nu_R^*}
             V(\one_A-\lambda_{A,\mathfrak b})\right|
       &\le\|V\|_\infty\mathcal E_n(L,\mathfrak b,J)
                                                       \label{eq:av-observation-error}
\end{align}
for every bounded Borel observation $V$ of the original trajectories.
All actual collision counts remain in $\mu_{n,R}$, $A$ and
$\lambda_{A,\mathfrak b}$.  No positive lower bound on any interval
length, incidence or clearance is required.
\end{theorem}
\begin{proof}
For the positive proof restriction $Q^L$, integrate the all-component
translation estimate in \eqref{eq:av-return-bv} against
$k_{\mathfrak b}$.  Tonelli and the triangle inequality give
\[
 \|Q^L-k_{\mathfrak b}*Q^L\|_{\rm TV}
 \le\mathcal V_L\int|h|k_{\mathfrak b}(h)\,dh
 =\kappa_{\rm ker}\mathcal V_L/\mathfrak b.
\]
The positive measures $H^L$ and $k_{\mathfrak b}*H^L$ both have
mass $\delta$; their difference has variation at most $2\delta$.
This proves \eqref{eq:av-law-error} for the full law.

For a union $I$ of at most $J$ intervals, translation of its
indicator satisfies
\[
 |I\mathbin\triangle(I+h)|\le2J|h|.
\]
It follows first for one interval by its two endpoints and then
for a finite union by subadditivity; it remains true for overlapping
interval presentations.  For a nonnegative BV component density
$p$, its compact physical support and Theorem~\ref{thm:av-all-bv} give
$\|p\|_\infty\le\|Dp\|_{\rm TV}/2$.  Therefore
\[
 \int p(t)|\one_I(t)-\one_I(t-h)|\,dt
             \le J|h|\|Dp\|_{\rm TV}.
\]
Positivity of the kernel permits integration of this inequality
in $h$ and summation over every component of $Q^L$.  The resulting
bound is $\kappa_{\rm ker}J\mathcal V_L/\mathfrak b$.
On the removed positive tail, both the indicator and likelihood
belong to $[0,1]$, so the absolute difference costs at most its
mass $\delta$.  This proves \eqref{eq:av-event-error} on all the
original trajectories.  Multiplication by $|V|\le\|V\|_\infty$
proves \eqref{eq:av-observation-error}.  It uses no regularity or
independence of the observation and no BV assertion for $V$.
Every exchange in this proof is of nonnegative integrals or an
absolutely bounded integral at the stated finite parameters.
\end{proof}

\subsection{An exact four-coordinate frequency formula}

Put
\[
 \widehat\mu_{n,R}(\xi,\upsilon,s)=
 \mathbb E_{\nu_R^*}
 e^{i(\xi\cdot K_{n,R}+\upsilon N_{n,R}+sS_{n,R})}.
\]
For one exact label $(k,m)$ and one bounded interval union $I$, write
$\lambda_{k,m,I,\mathfrak b}$ for
$\one_{\{K_{n,R}=k,N_{n,R}=m\}}
 (k_{\mathfrak b}*\one_I)(S_{n,R})$.
\begin{corollary}[The full finite-band transform uses the true law]
\label{cor:av-exact-fourier}
The exact identity is
\begin{align}
 \mathbb E_{\nu_R^*}\lambda_{k,m,I,\mathfrak b}
  ={}&(2\pi)^{-4}\int_{[-\pi,\pi]^3}\int_{-\mathfrak b}^{\mathfrak b}
      e^{-i(\xi\cdot k+\upsilon m)}
      \eta(s/\mathfrak b)\widehat{\one_I}(-s)
      \widehat\mu_{n,R}(\xi,\upsilon,s)
                           \,ds\,d\xi_1\,d\xi_2\,d\upsilon.
                                                       \label{eq:av-exact-fourier}
\end{align}
No torus frequency or arithmetic class is omitted.  A general event
in Theorem~\ref{thm:av-likelihood} is the positive sum of these
one-label likelihoods.  No interchange of an infinite complex label
series with the displayed frequency integral is required.
\end{corollary}
\begin{proof}
Character orthogonality on the full three-dimensional torus gives
the exact label indicator.  Since $\eta$ is even,
\[
 (k_{\mathfrak b}*\one_I)(t)
 =(2\pi)^{-1}\int_{-\mathfrak b}^{\mathfrak b}
       e^{ist}\eta(s/\mathfrak b)\widehat{\one_I}(-s)\,ds.
\]
Multiply the two identities and integrate on the actual probability
space.  The frequency domain is bounded and
$|\widehat{\one_I}(s)|\le|I|$, so Fubini is justified without a
collision cutoff.  This proves \eqref{eq:av-exact-fourier} with its
factor $(2\pi)^{-4}$.  The last assertion follows by monotone
summation of nonnegative trajectory likelihoods.  A discrete label
having zero physical mass contributes zero.  This does not assert
that smoothing has zero mass on every zero-mass roof window within
an attainable discrete class.
\end{proof}

\begin{corollary}[Arbitrary polynomial normalized accuracy]
\label{cor:av-polynomial}
Fix $P>0$, $\kappa\ge0$ and $J_0\ge1$.  If
$J_n\le J_0n^\kappa$, choose
\begin{equation}\label{eq:av-polynomial-scales}
 L_n=\lceil n^{P+3}\rceil,\qquad
 \mathfrak b_n=n^{2P+\kappa+5}.
\end{equation}
Uniformly in the original radius, all labels, all interval events
as above and all $\|V\|_\infty\le1$, the errors in
\eqref{eq:av-law-error}--\eqref{eq:av-observation-error}, multiplied
by $n^2$, are $O(n^{-P})$.  The constants depend only on the
displayed fixed parameters and the original physical family.
\end{corollary}
\begin{proof}
For $n\ge1$, $n^{P+3}\le L_n\le2n^{P+3}$.  Consequently
\[
 n^2\delta_{n,R}(L_n)\le c_*^{-1}n^{-P},\qquad
 \frac{n^2\kappa_{\rm ker}J_nD_{\rm all}^{1}L_n}{\mathfrak b_n}
        \le2\kappa_{\rm ker}J_0D_{\rm all}^{1}n^{-P}.
\]
For the law error omit $J_n$ and use $n^{-\kappa}\le1$.
These are substitutions into inequalities already valid for every
finite $L$ and every positive bandwidth.  They exchange no
geometric outer limit and impose no lower microscopic roof width.
\end{proof}

\begin{corollary}[Fixed collision count and the original $m^2$ scale]
\label{cor:av-fixed-collision}
At any fixed true collision count $m$, consider the entire positive
exact-label subprobability with components $p^w_{n,k,m,R}(t)\,dt$,
where $w=1$ or the original fixed $W^\varepsilon$.  The conclusions
\eqref{eq:av-law-error}--\eqref{eq:av-observation-error} hold with
$\delta=0$ and $\mathcal V_L=D_{\rm all}^w$, using the original
physical source weight $w\eta_R(T_R^mx)d\nu_R^*$.
If $J_m\le J_0m^\kappa$ and
$\mathfrak b_m=m^{P+\kappa+2}$, the absolute event and observation
errors multiplied by $m^2$ are $O(m^{-P})$.  There is no collision
restriction or renormalization in this assertion.
\end{corollary}
\begin{proof}
Apply Theorem~\ref{thm:av-all-bv} directly to the full fixed-count
label vector and repeat the positive kernel argument.  The source
is a subprobability and every test remains on that source; its mass
need not be replaced by one in any equality.  Multiplying its
event bound by $m^2$ gives
$\kappa_{\rm ker}J_0D_{\rm all}^wm^{-P}$.
\end{proof}

\begin{corollary}[The actual conditional denominator]
\label{cor:av-denominator}
Let $H$ be any Borel information event on the original trajectories,
$a=\one_H\one_A$, $b=\one_H\lambda_{A,\mathfrak b}$, and
$p=\mathbb E a$, $p_{\mathfrak b}=\mathbb E b$.  If the applicable
absolute error bound is $E$, then
\[
 |p-p_{\mathfrak b}|\le E.
\]
If the actual selected mass satisfies $p>E$, both normalized
measures exist and their full variation distance obeys
\begin{equation}\label{eq:av-relative}
 \left\|\frac{a\,d\nu_R^*}{p}
       -\frac{b\,d\nu_R^*}{p_{\mathfrak b}}\right\|_{\rm TV}
                                  \le\frac{2E}{p}.
\end{equation}
If instead $p_{\mathfrak b}>E$ is verified, the bound
$2E/(p_{\mathfrak b}-E)$ is valid.  The same statements hold with
the fixed-collision source of Corollary~\ref{cor:av-fixed-collision}.
\end{corollary}
\begin{proof}
The positive trajectory estimate gives $\mathbb E|a-b|\le E$,
hence the mass bound.  If $p>E$, then $p_{\mathfrak b}\ge p-E>0$.
Add and subtract $b\,d\nu_R^*/p$.  The first variation cost is
at most $E/p$; the second is
$p_{\mathfrak b}|p^{-1}-p_{\mathfrak b}^{-1}|
=|p-p_{\mathfrak b}|/p\le E/p$.  This proves
\eqref{eq:av-relative}.  If $p_{\mathfrak b}>E$, use
$p\ge p_{\mathfrak b}-E>0$ in the same calculation.  The proof
requires only the stated positive base measure, so applies also to
the fixed-collision case.
\end{proof}

\paragraph{Relation to the original raw target.}
The interval event, each collision and its unit-speed physical time
remain exact throughout.  Convolution defines an analytical
likelihood on those trajectories; it is not a new physical clock or
a changed return condition.  The absolute estimates are uniform even
when the original roof intervals shrink arbitrarily.  Dividing by
their lengths, or by a small selected-event probability, requires
the corresponding additional estimate; neither is implicit here.
In particular each information restriction has its own actual
denominator.  A lower bound for a larger event cannot be inherited
by restriction, and a zero arithmetic event cannot be conditioned
upon.  These conclusions concern $\nu_R^*$; they do not automatically
transfer to a stationary entrance law with an unbounded return-roof
density.  The full joint spectral evaluation of
\eqref{eq:av-exact-fourier}, every relevant arithmetic peak, the
original endpoint-layer shortage and the four-coordinate raw local
inverse remain separate requirements.  No averaged limit or finite
frequency calculation is being substituted for their proofs.
